% GENERATED FILE. Edit canonical lesson frontmatter, then run npm run generate:books.
% canonical-entries: 109

\chapter*{Glossary}
\addcontentsline{toc}{chapter}{Glossary}
\markboth{Glossary}{Glossary}

This glossary collects every word and phrase taught in the book. Pronunciation appears when it differs from the written form; chapter numbers show where each entry is introduced.

\noindent\begin{minipage}[t]{\linewidth}
\raggedright
\textbf{\hi{अच्छा}}\enspace\emph{acchā}\par
\small good\par
\footnotesize Introduced in Chapter 41.
\end{minipage}\par\medskip

\noindent\begin{minipage}[t]{\linewidth}
\raggedright
\textbf{\hi{आना}}\enspace\emph{ānā}\par
\small to come\par
\footnotesize Introduced in Chapter 42.
\end{minipage}\par\medskip

\noindent\begin{minipage}[t]{\linewidth}
\raggedright
\textbf{\hi{बच्चा}}\enspace\emph{bacca}\par
\small child — masculine, with a feminine \hi{बच्ची}, and a Persian word that landed on top of an inherited Sanskrit cousin\par
\footnotesize Introduced in Chapter 39.
\end{minipage}\par\medskip

\noindent\begin{minipage}[t]{\linewidth}
\raggedright
\textbf{\hi{बैठना}}\enspace\emph{baiṭhnā}\par
\small to sit\par
\footnotesize Introduced in Chapter 42.
\end{minipage}\par\medskip

\noindent\begin{minipage}[t]{\linewidth}
\raggedright
\textbf{\hi{बड़ा}}\enspace\emph{baṛā}\par
\small big\par
\footnotesize Introduced in Chapter 41.
\end{minipage}\par\medskip

\noindent\begin{minipage}[t]{\linewidth}
\raggedright
\textbf{\hi{बेटा}}\enspace\emph{beṭā}\par
\small son\par
\footnotesize Introduced in Chapter 44.
\end{minipage}\par\medskip

\noindent\begin{minipage}[t]{\linewidth}
\raggedright
\textbf{\hi{बेटी}}\enspace\emph{beṭī}\par
\small daughter\par
\footnotesize Introduced in Chapter 44.
\end{minipage}\par\medskip

\noindent\begin{minipage}[t]{\linewidth}
\raggedright
\textbf{\hi{बोलना}}\enspace\emph{bolnā}\par
\small to speak\par
\footnotesize Introduced in Chapter 43.
\end{minipage}\par\medskip

\noindent\begin{minipage}[t]{\linewidth}
\raggedright
\textbf{\hi{छोटा}}\enspace\emph{choṭā}\par
\small small\par
\footnotesize Introduced in Chapter 41.
\end{minipage}\par\medskip

\noindent\begin{minipage}[t]{\linewidth}
\raggedright
\textbf{\hi{दाँत}}\enspace\emph{danta}\par
\small tooth — masculine, though it is built exactly like the feminine \hi{आँख}, and a cousin of Latin dēns and English tooth\par
\footnotesize Introduced in Chapter 38.
\end{minipage}\par\medskip

\noindent\begin{minipage}[t]{\linewidth}
\raggedright
\textbf{\hi{देखना}}\enspace\emph{dekhnā}\par
\small to look, to see\par
\footnotesize Introduced in Chapter 43.
\end{minipage}\par\medskip

\noindent\begin{minipage}[t]{\linewidth}
\raggedright
\textbf{\hi{देना}}\enspace\emph{denā}\par
\small to give\par
\footnotesize Introduced in Chapter 42.
\end{minipage}\par\medskip

\noindent\begin{minipage}[t]{\linewidth}
\raggedright
\textbf{\hi{घंटा}}\enspace\emph{ghaṇṭā}\par
\small hour — literally \textquotedblleft{}bell,\textquotedblright{} from the historical practice of striking a bell to mark the hour\par
\footnotesize Introduced in Chapter 18.
\end{minipage}\par\medskip

\noindent\begin{minipage}[t]{\linewidth}
\raggedright
\textbf{\hi{है}}\enspace\emph{hai}\par
\small is\par
\footnotesize Introduced in Chapter 2.
\end{minipage}\par\medskip

\noindent\begin{minipage}[t]{\linewidth}
\raggedright
\textbf{\hi{हरा}}\enspace\emph{harā}\par
\small green — a deep cousin of German gelb and English yellow from an ancient shining-color root\par
\footnotesize Introduced in Chapter 24.
\end{minipage}\par\medskip

\noindent\begin{minipage}[t]{\linewidth}
\raggedright
\textbf{\hi{हूँ}}\enspace\emph{hūṁ}\par
\small am (1st-person of \hi{होना}, \textquotedblleft{}to be\textquotedblright{})\par
\footnotesize Introduced in Chapter 3.
\end{minipage}\par\medskip

\noindent\begin{minipage}[t]{\linewidth}
\raggedright
\textbf{\hi{कहाँ}}\enspace\emph{kahā̃}\par
\small where? — asking about a place\par
\footnotesize Introduced in Chapter 40.
\end{minipage}\par\medskip

\noindent\begin{minipage}[t]{\linewidth}
\raggedright
\textbf{\hi{कैसे}}\enspace\emph{kaise}\par
\small how (also \hi{कैसा}/\hi{कैसी}, agreeing)\par
\footnotesize Introduced in Chapter 3.
\end{minipage}\par\medskip

\noindent\begin{minipage}[t]{\linewidth}
\raggedright
\textbf{\hi{कौन}}\enspace\emph{kaun}\par
\small who? — asking about a person\par
\footnotesize Introduced in Chapter 40.
\end{minipage}\par\medskip

\noindent\begin{minipage}[t]{\linewidth}
\raggedright
\textbf{\hi{खाना}}\enspace\emph{khānā}\par
\small food — masculine, from a Sanskrit verb of eating, and not to be confused with its Persian look-alike meaning \textquotedblleft{}house\textquotedblright{}\par
\footnotesize Introduced in Chapter 37.
\end{minipage}\par\medskip

\noindent\begin{minipage}[t]{\linewidth}
\raggedright
\textbf{\hi{खड़ा} \hi{होना}}\enspace\emph{khaṛā honā}\par
\small to stand\par
\footnotesize Introduced in Chapter 42.
\end{minipage}\par\medskip

\noindent\begin{minipage}[t]{\linewidth}
\raggedright
\textbf{\hi{कृपया}}\enspace\emph{kṛpayā}\par
\small please (kṛpayā — \textquotedblleft{}kindly,\textquotedblright{} from kṛpā \textquotedblleft{}grace/compassion\textquotedblright{})\par
\footnotesize Introduced in Chapter 8.
\end{minipage}\par\medskip

\noindent\begin{minipage}[t]{\linewidth}
\raggedright
\textbf{\hi{कुर्सी}}\enspace\emph{kursī}\par
\small chair — your first feminine noun, and a word that has been passing between Middle Eastern languages for four thousand years\par
\footnotesize Introduced in Chapter 36.
\end{minipage}\par\medskip

\noindent\begin{minipage}[t]{\linewidth}
\raggedright
\textbf{\hi{कुत्ता}}\enspace\emph{kuttā}\par
\small dog — a third uncertain-origin everyday word displacing an inherited PIE dog-word\par
\footnotesize Introduced in Chapter 23.
\end{minipage}\par\medskip

\noindent\begin{minipage}[t]{\linewidth}
\raggedright
\textbf{\hi{क्या}}\enspace\emph{kyā}\par
\small what\par
\footnotesize Introduced in Chapter 2.
\end{minipage}\par\medskip

\noindent\begin{minipage}[t]{\linewidth}
\raggedright
\textbf{\hi{लेना}}\enspace\emph{lenā}\par
\small to take — worn down from Sanskrit labh-, which also sits, undisguised, in the noun \hi{लाभ}\par
\footnotesize Introduced in Chapter 35.
\end{minipage}\par\medskip

\noindent\begin{minipage}[t]{\linewidth}
\raggedright
\textbf{\hi{लिखना}}\enspace\emph{likhnā}\par
\small to write\par
\footnotesize Introduced in Chapter 43.
\end{minipage}\par\medskip

\noindent\begin{minipage}[t]{\linewidth}
\raggedright
\textbf{\hi{लोग}}\enspace\emph{log}\par
\small people\par
\footnotesize Introduced in Chapter 44.
\end{minipage}\par\medskip

\noindent\begin{minipage}[t]{\linewidth}
\raggedright
\textbf{\hi{मैं}}\enspace\emph{maiṁ}\par
\small I\par
\footnotesize Introduced in Chapter 3.
\end{minipage}\par\medskip

\noindent\begin{minipage}[t]{\linewidth}
\raggedright
\textbf{\hi{मेरा}}\enspace\emph{merā}\par
\small my\par
\footnotesize Introduced in Chapter 2.
\end{minipage}\par\medskip

\noindent\begin{minipage}[t]{\linewidth}
\raggedright
\textbf{\hi{मिलेंगे}}\enspace\emph{milenge}\par
\small (we) will meet\par
\footnotesize Introduced in Chapter 4.
\end{minipage}\par\medskip

\noindent\begin{minipage}[t]{\linewidth}
\raggedright
\textbf{\hi{नमस्ते}}\enspace\emph{namaste}\par
\small hello / goodbye (namaste — \textquotedblleft{}I bow to you\textquotedblright{})\par
\footnotesize Introduced in Chapter 1.
\end{minipage}\par\medskip

\noindent\begin{minipage}[t]{\linewidth}
\raggedright
\textbf{\hi{नया}}\enspace\emph{nayā}\par
\small new\par
\footnotesize Introduced in Chapter 41.
\end{minipage}\par\medskip

\noindent\begin{minipage}[t]{\linewidth}
\raggedright
\textbf{\hi{पानी} \hi{रोटी}}\enspace\emph{pānī roṭī}\par
\small water (literally \textquotedblleft{}the drinkable one,\textquotedblright{} not the ancient Sanskrit word for water) and bread/flatbread\par
\footnotesize Introduced in Chapter 15.
\end{minipage}\par\medskip

\noindent\begin{minipage}[t]{\linewidth}
\raggedright
\textbf{\hi{परिवार}}\enspace\emph{parivār}\par
\small family\par
\footnotesize Introduced in Chapter 44.
\end{minipage}\par\medskip

\noindent\begin{minipage}[t]{\linewidth}
\raggedright
\textbf{\hi{पुराना}}\enspace\emph{purānā}\par
\small old — of a thing, not of a person\par
\footnotesize Introduced in Chapter 41.
\end{minipage}\par\medskip

\noindent\begin{minipage}[t]{\linewidth}
\raggedright
\textbf{\hi{सुनना}}\enspace\emph{sunnā}\par
\small to listen, to hear\par
\footnotesize Introduced in Chapter 43.
\end{minipage}\par\medskip

\noindent\begin{minipage}[t]{\linewidth}
\raggedright
\textbf{\hi{वह}}\enspace\emph{vah}\par
\small that one — the thing over there\par
\footnotesize Introduced in Chapter 40.
\end{minipage}\par\medskip

\noindent\begin{minipage}[t]{\linewidth}
\raggedright
\textbf{\hi{वहाँ}}\enspace\emph{vahā̃}\par
\small there — where I am not\par
\footnotesize Introduced in Chapter 40.
\end{minipage}\par\medskip

\noindent\begin{minipage}[t]{\linewidth}
\raggedright
\textbf{\hi{यह}}\enspace\emph{yah}\par
\small this one — the thing near me\par
\footnotesize Introduced in Chapter 40.
\end{minipage}\par\medskip

\noindent\begin{minipage}[t]{\linewidth}
\raggedright
\textbf{\hi{यहाँ}}\enspace\emph{yahā̃}\par
\small here — where I am\par
\footnotesize Introduced in Chapter 40.
\end{minipage}\par\medskip

\noindent\begin{minipage}[t]{\linewidth}
\raggedright
\textbf{\hi{अलविदा}}\par
\small goodbye (alvidā — a final farewell, Persian/Arabic-derived)\par
\footnotesize Introduced in Chapter 1.
\end{minipage}\par\medskip

\noindent\begin{minipage}[t]{\linewidth}
\raggedright
\textbf{\hi{आँख}}\par
\small eye — feminine, and a word-for-word cousin of English \textquotedblleft{}eye\textquotedblright{} and Latin \textquotedblleft{}oculus\textquotedblright{}\par
\footnotesize Introduced in Chapter 38.
\end{minipage}\par\medskip

\noindent\begin{minipage}[t]{\linewidth}
\raggedright
\textbf{\hi{आदमी}}\par
\small man, person — masculine despite its -ī ending, and literally \textquotedblleft{}of Adam\textquotedblright{}\par
\footnotesize Introduced in Chapter 39.
\end{minipage}\par\medskip

\noindent\begin{minipage}[t]{\linewidth}
\raggedright
\textbf{\hi{आप} / \hi{तुम}}\par
\small you (respectful / familiar)\par
\footnotesize Introduced in Chapter 2.
\end{minipage}\par\medskip

\noindent\begin{minipage}[t]{\linewidth}
\raggedright
\textbf{\hi{आप} \hi{कैसे} \hi{हैं}?}\par
\small how are you? (respectful)\par
\footnotesize Introduced in Chapter 3.
\end{minipage}\par\medskip

\noindent\begin{minipage}[t]{\linewidth}
\raggedright
\textbf{\hi{आपका} \hi{नाम} \hi{क्या} \hi{है}?}\par
\small what's your name?\par
\footnotesize Introduced in Chapter 2.
\end{minipage}\par\medskip

\noindent\begin{minipage}[t]{\linewidth}
\raggedright
\textbf{\hi{आपका} \hi{स्वागत} \hi{है}}\par
\small you're welcome (also \textquotedblleft{}welcome!\textquotedblright{})\par
\footnotesize Introduced in Chapter 3.
\end{minipage}\par\medskip

\noindent\begin{minipage}[t]{\linewidth}
\raggedright
\textbf{\hi{उम्र}}\par
\small age — the same word as Arabic's ʿumr, borrowed through Persian, beside formal Sanskrit āyu\par
\footnotesize Introduced in Chapter 19.
\end{minipage}\par\medskip

\noindent\begin{minipage}[t]{\linewidth}
\raggedright
\textbf{\hi{एक} \hi{दो} \hi{तीन} \hi{चार} \hi{पाँच}}\par
\small one to five, with familiar Devanagari pieces and the chandrabindu in \hi{पाँच}\par
\footnotesize Introduced in Chapter 6.
\end{minipage}\par\medskip

\noindent\begin{minipage}[t]{\linewidth}
\raggedright
\textbf{\hi{औरत}}\par
\small woman — feminine, an Arabic word that only came to mean \textquotedblleft{}woman\textquotedblright{} in Persian, and one of three Hindi words separated by register\par
\footnotesize Introduced in Chapter 39.
\end{minipage}\par\medskip

\noindent\begin{minipage}[t]{\linewidth}
\raggedright
\textbf{\hi{किताब}}\par
\small book — feminine, built on the Arabic root for writing, and one of three Hindi words for a book that differ by register\par
\footnotesize Introduced in Chapter 37.
\end{minipage}\par\medskip

\noindent\begin{minipage}[t]{\linewidth}
\raggedright
\textbf{\hi{कमरा}}\par
\small room — masculine, and the same Latin word as English \textquotedblleft{}camera\textquotedblright{}, brought to India by Portuguese ships\par
\footnotesize Introduced in Chapter 36.
\end{minipage}\par\medskip

\noindent\begin{minipage}[t]{\linewidth}
\raggedright
\textbf{\hi{करना}}\par
\small to do, to make (and \hi{काम} \hi{करना}, \textquotedblleft{}to work\textquotedblright{})\par
\footnotesize Introduced in Chapter 5.
\end{minipage}\par\medskip

\noindent\begin{minipage}[t]{\linewidth}
\raggedright
\textbf{\hi{कल} \hi{मिलते} \hi{हैं}}\par
\small see you tomorrow (and the \hi{कल} puzzle)\par
\footnotesize Introduced in Chapter 4.
\end{minipage}\par\medskip

\noindent\begin{minipage}[t]{\linewidth}
\raggedright
\textbf{\hi{काला} \hi{सफ़ेद}}\par
\small black and white — one Sanskrit word tangled up with time and death, one a Persian loan\par
\footnotesize Introduced in Chapter 11.
\end{minipage}\par\medskip

\noindent\begin{minipage}[t]{\linewidth}
\raggedright
\textbf{\hi{ख़ुशी}}\par
\small happiness / pleased to meet you\par
\footnotesize Introduced in Chapter 2.
\end{minipage}\par\medskip

\noindent\begin{minipage}[t]{\linewidth}
\raggedright
\textbf{\hi{ग्यारह} — \hi{बीस}}\par
\small 11-20 — genuinely irregular, must be memorized individually; but \hi{उन्नीस} (19) most likely preserves Sanskrit's OWN subtractive naming, the same logic as Latin's ūndēvīgintī\par
\footnotesize Introduced in Chapter 22.
\end{minipage}\par\medskip

\noindent\begin{minipage}[t]{\linewidth}
\raggedright
\textbf{\hi{घर}}\par
\small house, home — masculine, and from a Sanskrit word that meant a fenced enclosure before it meant a building\par
\footnotesize Introduced in Chapter 36.
\end{minipage}\par\medskip

\noindent\begin{minipage}[t]{\linewidth}
\raggedright
\textbf{\hi{चाय}}\par
\small tea — feminine, and one of the two names the world uses for the same leaf\par
\footnotesize Introduced in Chapter 37.
\end{minipage}\par\medskip

\noindent\begin{minipage}[t]{\linewidth}
\raggedright
\textbf{\hi{चलता} \hi{हूँ}}\par
\small I'll be off (lit. \textquotedblleft{}I go\textquotedblright{}), gendered\par
\footnotesize Introduced in Chapter 4.
\end{minipage}\par\medskip

\noindent\begin{minipage}[t]{\linewidth}
\raggedright
\textbf{\hi{छह} \hi{सात} \hi{आठ} \hi{नौ} \hi{दस}}\par
\small six to ten, with seven and eight repeating the same Prakrit weight trade as three\par
\footnotesize Introduced in Chapter 21.
\end{minipage}\par\medskip

\noindent\begin{minipage}[t]{\linewidth}
\raggedright
\textbf{\hi{जनवरी} \hi{फ़रवरी} \hi{मार्च} \hi{अप्रैल} \hi{मई} \hi{जून} \hi{जुलाई} \hi{अगस्त} \hi{सितंबर} \hi{अक्टूबर} \hi{नवंबर} \hi{दिसंबर}}\par
\small the Gregorian months — borrowed English/international names, running ALONGSIDE the traditional Vikram Samvat calendar's own six-ritu-linked months\par
\footnotesize Introduced in Chapter 16.
\end{minipage}\par\medskip

\noindent\begin{minipage}[t]{\linewidth}
\raggedright
\textbf{\hi{ठीक}}\par
\small fine, well, correct — and the reply \textquotedblleft{}\hi{मैं} \hi{ठीक} \hi{हूँ}\textquotedblright{}\par
\footnotesize Introduced in Chapter 3.
\end{minipage}\par\medskip

\noindent\begin{minipage}[t]{\linewidth}
\raggedright
\textbf{\hi{दूध}}\par
\small milk — masculine, and literally \textquotedblleft{}the milked thing\textquotedblright{}, a past participle that hardened into a noun\par
\footnotesize Introduced in Chapter 37.
\end{minipage}\par\medskip

\noindent\begin{minipage}[t]{\linewidth}
\raggedright
\textbf{\hi{दिन}}\par
\small \textquotedblleft{}day\textquotedblright{} — shares its ultimate PIE root with Latin's diēs, but through a separate derivational branch (din patterns with Baltic/Slavic \textquotedblleft{}day\textquotedblright{}-words); Sanskrit's own \hi{दिव्}/\hi{द्यु} are actually the closer, more direct cousins of diēs\par
\footnotesize Introduced in Chapter 25.
\end{minipage}\par\medskip

\noindent\begin{minipage}[t]{\linewidth}
\raggedright
\textbf{\hi{दोपहर}}\par
\small \textquotedblleft{}afternoon\textquotedblright{} — the same word already met meaning precisely \textquotedblleft{}noon\textquotedblright{} (two pahars after sunrise, HI-C17), now widened in modern everyday usage to cover the whole afternoon stretch, not just the exact midday moment\par
\footnotesize Introduced in Chapter 30.
\end{minipage}\par\medskip

\noindent\begin{minipage}[t]{\linewidth}
\raggedright
\textbf{\hi{दोपहर} \hi{आधी} \hi{रात}}\par
\small noon (literally \textquotedblleft{}two pahars,\textquotedblright{} a traditional \textasciitilde{}3-hour time unit) and midnight (\textquotedblleft{}half night,\textquotedblright{} fully transparent)\par
\footnotesize Introduced in Chapter 17.
\end{minipage}\par\medskip

\noindent\begin{minipage}[t]{\linewidth}
\raggedright
\textbf{\hi{दरवाज़ा}}\par
\small door — masculine, a Persian loan whose first syllable is the same ancient word as English \textquotedblleft{}door\textquotedblright{}\par
\footnotesize Introduced in Chapter 36.
\end{minipage}\par\medskip

\noindent\begin{minipage}[t]{\linewidth}
\raggedright
\textbf{\hi{दोस्त}}\par
\small friend — a Persian word for \textquotedblleft{}the one chosen\textquotedblright{}, masculine by default but taking either possessive depending on who the friend is\par
\footnotesize Introduced in Chapter 39.
\end{minipage}\par\medskip

\noindent\begin{minipage}[t]{\linewidth}
\raggedright
\textbf{\hi{धन्यवाद}}\par
\small thank you (dhanyavād — formal, Sanskritic)\par
\footnotesize Introduced in Chapter 1.
\end{minipage}\par\medskip

\noindent\begin{minipage}[t]{\linewidth}
\raggedright
\textbf{\hi{नाम}}\par
\small name\par
\footnotesize Introduced in Chapter 2.
\end{minipage}\par\medskip

\noindent\begin{minipage}[t]{\linewidth}
\raggedright
\textbf{\hi{नमस्कार}}\par
\small hello (namaskār — more formal)\par
\footnotesize Introduced in Chapter 1.
\end{minipage}\par\medskip

\noindent\begin{minipage}[t]{\linewidth}
\raggedright
\textbf{\hi{नहीं}}\par
\small no / not (nahīṃ)\par
\footnotesize Introduced in Chapter 7.
\end{minipage}\par\medskip

\noindent\begin{minipage}[t]{\linewidth}
\raggedright
\textbf{\hi{पूछना}}\par
\small to ask — the same inherited verb as English \textquotedblleft{}pray\textquotedblright{}, which once meant nothing more than asking\par
\footnotesize Introduced in Chapter 35.
\end{minipage}\par\medskip

\noindent\begin{minipage}[t]{\linewidth}
\raggedright
\textbf{\hi{पेट}}\par
\small stomach, belly — masculine, and the one word in this chapter with no clean ancestry at all\par
\footnotesize Introduced in Chapter 38.
\end{minipage}\par\medskip

\noindent\begin{minipage}[t]{\linewidth}
\raggedright
\textbf{\hi{पढ़ना}}\par
\small to read — and equally \textquotedblleft{}to study\textquotedblright{}, because the Sanskrit verb underneath it meant \textquotedblleft{}to recite aloud\textquotedblright{}\par
\footnotesize Introduced in Chapter 34.
\end{minipage}\par\medskip

\noindent\begin{minipage}[t]{\linewidth}
\raggedright
\textbf{\hi{पिता} \hi{माता}}\par
\small father and mother — Sanskrit cousins of Latin pater/mater and English father/mother\par
\footnotesize Introduced in Chapter 12.
\end{minipage}\par\medskip

\noindent\begin{minipage}[t]{\linewidth}
\raggedright
\textbf{\hi{पैर}}\par
\small foot — masculine, and NOT the descendant of Sanskrit pāda; that honour belongs to \hi{पाँव}\par
\footnotesize Introduced in Chapter 38.
\end{minipage}\par\medskip

\noindent\begin{minipage}[t]{\linewidth}
\raggedright
\textbf{\hi{पसंद}}\par
\small liked, pleasing — a Persian noun, not a verb, in a sentence where the thing you like is the subject and you are only the person it happens to\par
\footnotesize Introduced in Chapter 35.
\end{minipage}\par\medskip

\noindent\begin{minipage}[t]{\linewidth}
\raggedright
\textbf{\hi{फिर}}\par
\small again, then\par
\footnotesize Introduced in Chapter 4.
\end{minipage}\par\medskip

\noindent\begin{minipage}[t]{\linewidth}
\raggedright
\textbf{\hi{फिर} \hi{मिलेंगे}}\par
\small we'll meet again (goodbye)\par
\footnotesize Introduced in Chapter 4.
\end{minipage}\par\medskip

\noindent\begin{minipage}[t]{\linewidth}
\raggedright
\textbf{\hi{बोलता} \hi{हूँ}, \hi{बोलती} \hi{हूँ}, \hi{बोलते} \hi{हैं}}\par
\small the present-habitual template — stem + -tā/-tī/-te + honā\par
\footnotesize Introduced in Chapter 5.
\end{minipage}\par\medskip

\noindent\begin{minipage}[t]{\linewidth}
\raggedright
\textbf{\hi{बोलना}}\par
\small to speak\par
\footnotesize Introduced in Chapter 5.
\end{minipage}\par\medskip

\noindent\begin{minipage}[t]{\linewidth}
\raggedright
\textbf{\hi{भाई} \hi{बहन}}\par
\small brother (a PIE cousin of English \textquotedblleft{}brother\textquotedblright{}) and sister (a DIFFERENT Sanskrit word entirely, not the PIE cousin of \textquotedblleft{}sister\textquotedblright{})\par
\footnotesize Introduced in Chapter 12.
\end{minipage}\par\medskip

\noindent\begin{minipage}[t]{\linewidth}
\raggedright
\textbf{\hi{मैं} \hi{हिंदी} \hi{बोलता} \hi{हूँ}}\par
\small I speak Hindi (gendered)\par
\footnotesize Introduced in Chapter 5.
\end{minipage}\par\medskip

\noindent\begin{minipage}[t]{\linewidth}
\raggedright
\textbf{\hi{मदद} \hi{करना}}\par
\small to help — a noun plus \hi{करना}, which is how Hindi turns almost anything into a verb\par
\footnotesize Introduced in Chapter 35.
\end{minipage}\par\medskip

\noindent\begin{minipage}[t]{\linewidth}
\raggedright
\textbf{\hi{माफ़} \hi{कीजिए}}\par
\small please forgive me / sorry (māf kījiye — \textquotedblleft{}please do forgiveness\textquotedblright{})\par
\footnotesize Introduced in Chapter 9.
\end{minipage}\par\medskip

\noindent\begin{minipage}[t]{\linewidth}
\raggedright
\textbf{\hi{मेरा} \hi{नाम} … \hi{है}}\par
\small my name is…\par
\footnotesize Introduced in Chapter 2.
\end{minipage}\par\medskip

\noindent\begin{minipage}[t]{\linewidth}
\raggedright
\textbf{\hi{मौसम}}\par
\small weather — the SAME root as English \textquotedblleft{}monsoon\textquotedblright{}; rain has TWO words that sound alike but are NOT related, a real false-friend trap\par
\footnotesize Introduced in Chapter 20.
\end{minipage}\par\medskip

\noindent\begin{minipage}[t]{\linewidth}
\raggedright
\textbf{\hi{रात}}\par
\small \textquotedblleft{}night\textquotedblright{} — already met inside aadhi raat (\textquotedblleft{}midnight\textquotedblright{}); despite meaning the same thing as Latin's nox, in a related Indo-European language, it comes from a COMPLETELY DIFFERENT PIE root — a genuine false cousin\par
\footnotesize Introduced in Chapter 26.
\end{minipage}\par\medskip

\noindent\begin{minipage}[t]{\linewidth}
\raggedright
\textbf{\hi{रहना}}\par
\small to live, to stay\par
\footnotesize Introduced in Chapter 5.
\end{minipage}\par\medskip

\noindent\begin{minipage}[t]{\linewidth}
\raggedright
\textbf{\hi{लिखना}}\par
\small to write — from a Sanskrit verb meaning \textquotedblleft{}to scratch\textquotedblright{}, the same picture five language families reached independently\par
\footnotesize Introduced in Chapter 34.
\end{minipage}\par\medskip

\noindent\begin{minipage}[t]{\linewidth}
\raggedright
\textbf{\hi{लाल} \hi{नीला}}\par
\small red and blue — a Persian gemstone-name, and the Sanskrit word behind the world's word for \textquotedblleft{}indigo\textquotedblright{}\par
\footnotesize Introduced in Chapter 11.
\end{minipage}\par\medskip

\noindent\begin{minipage}[t]{\linewidth}
\raggedright
\textbf{\hi{वसंत} \hi{ग्रीष्म} \hi{वर्षा} \hi{शरद्} \hi{हेमंत} \hi{शिशिर}}\par
\small SIX traditional seasons, not four — India's own \hi{ऋतु} system, which the Western four-season frame doesn't map onto cleanly\par
\footnotesize Introduced in Chapter 14.
\end{minipage}\par\medskip

\noindent\begin{minipage}[t]{\linewidth}
\raggedright
\textbf{\hi{शुक्रिया}}\par
\small thanks (shukriyā — everyday, Persian/Arabic-derived)\par
\footnotesize Introduced in Chapter 1.
\end{minipage}\par\medskip

\noindent\begin{minipage}[t]{\linewidth}
\raggedright
\textbf{\hi{शनिवार} \hi{रविवार}}\par
\small Saturday and Sunday — completing Hindi's planet-week\par
\footnotesize Introduced in Chapter 10.
\end{minipage}\par\medskip

\noindent\begin{minipage}[t]{\linewidth}
\raggedright
\textbf{\hi{शुभ} \hi{दोपहर}}\par
\small \textquotedblleft{}good afternoon\textquotedblright{} (shubh dopahar) — uses \hi{दोपहर} in its modern, widened \textquotedblleft{}afternoon\textquotedblright{} sense (HI-C30), not the narrow etymological \textquotedblleft{}noon\textquotedblright{}; the rarest of this arc's three \hi{शुभ}+noun greetings in actual casual use, with \hi{नमस्ते}/\hi{नमस्कार} covering the afternoon in practice\par
\footnotesize Introduced in Chapter 33.
\end{minipage}\par\medskip

\noindent\begin{minipage}[t]{\linewidth}
\raggedright
\textbf{\hi{शुभ} \hi{रात्रि}}\par
\small good night" — literally \textquotedblleft{}auspicious night,\textquotedblright{} using the formal literary raatri, not the everyday raat just taught, plus shubh, an \textquotedblleft{}auspicious\textquotedblright{} word that started life meaning "shining\par
\footnotesize Introduced in Chapter 27.
\end{minipage}\par\medskip

\noindent\begin{minipage}[t]{\linewidth}
\raggedright
\textbf{\hi{शुभ} \hi{संध्या}}\par
\small \textquotedblleft{}good evening\textquotedblright{} — sandhyā is the Sanskrit junction of day and night, with the register-crossing doublet sāñjh\par
\footnotesize Introduced in Chapter 32.
\end{minipage}\par\medskip

\noindent\begin{minipage}[t]{\linewidth}
\raggedright
\textbf{\hi{शाम}}\par
\small \textquotedblleft{}evening\textquotedblright{} (shaam) — another genuine Persian loanword, like subah; looks tempting to connect to Sanskrit \hi{श्यामा} (śyāmā, \textquotedblleft{}night/dark\textquotedblright{}), but that's a false cognate — different root entirely, despite the resemblance\par
\footnotesize Introduced in Chapter 29.
\end{minipage}\par\medskip

\noindent\begin{minipage}[t]{\linewidth}
\raggedright
\textbf{\hi{सोचना}}\par
\small to think — a verb that began life meaning \textquotedblleft{}to grieve,\textquotedblright{} and still keeps the worry inside the noun \hi{सोच}\par
\footnotesize Introduced in Chapter 34.
\end{minipage}\par\medskip

\noindent\begin{minipage}[t]{\linewidth}
\raggedright
\textbf{\hi{सुप्रभात}}\par
\small \textquotedblleft{}good morning\textquotedblright{} (suprabhat) — \hi{सु} (\textquotedblleft{}good\textquotedblright{}) + \hi{प्रभात} (prabhāt, native Sanskrit \textquotedblleft{}dawn,\textquotedblright{} from prabhā, \textquotedblleft{}light, splendour,\textquotedblright{} itself from bhā, \textquotedblleft{}to shine\textquotedblright{}) — a completely different word from \hi{सुबह}, the Persian-loan everyday noun; skews formal/written (radio, greeting cards) rather than casual spoken Hindi, where \hi{नमस्ते} covers mornings too — but there's no strong evidence it's being displaced by English the way some other languages' night greetings were found to be\par
\footnotesize Introduced in Chapter 31.
\end{minipage}\par\medskip

\noindent\begin{minipage}[t]{\linewidth}
\raggedright
\textbf{\hi{सुबह}}\par
\small \textquotedblleft{}morning\textquotedblright{} (subah) — a Persian/Arabic loanword, NOT Sanskrit tatsama like din/raat; shares the Arabic root \ar{ص}-\ar{ب}-\ar{ح} with ṣabāḥ (\textquotedblleft{}morning\textquotedblright{}) and with the tuṣbiḥ inside the Arabic \textquotedblleft{}good night\textquotedblright{} phrase\par
\footnotesize Introduced in Chapter 28.
\end{minipage}\par\medskip

\noindent\begin{minipage}[t]{\linewidth}
\raggedright
\textbf{\hi{समझना}}\par
\small to understand — literally \textquotedblleft{}to wake up fully,\textquotedblright{} on the same root as Buddha\par
\footnotesize Introduced in Chapter 34.
\end{minipage}\par\medskip

\noindent\begin{minipage}[t]{\linewidth}
\raggedright
\textbf{\hi{सोमवार} \hi{मंगलवार} \hi{बुधवार} \hi{गुरुवार} \hi{शुक्रवार}}\par
\small Monday–Friday — Hindi's own planet-week, independent of Rome's\par
\footnotesize Introduced in Chapter 10.
\end{minipage}\par\medskip

\noindent\begin{minipage}[t]{\linewidth}
\raggedright
\textbf{\hi{सिर}}\par
\small the head — from Sanskrit śiras, worn down through a thousand years of everyday speech\par
\footnotesize Introduced in Chapter 13.
\end{minipage}\par\medskip

\noindent\begin{minipage}[t]{\linewidth}
\raggedright
\textbf{\hi{हाँ}}\par
\small yes (hāṃ)\par
\footnotesize Introduced in Chapter 7.
\end{minipage}\par\medskip

\noindent\begin{minipage}[t]{\linewidth}
\raggedright
\textbf{\hi{हाथ}}\par
\small the hand — from Sanskrit hasta, with a real (if non-obvious) sound-change story\par
\footnotesize Introduced in Chapter 13.
\end{minipage}\par\medskip
